\section{Ordinary labels, single cycles, and block cores}\label{sec:labels}

Throughout this section $p$ is odd and different from the defining
characteristic.  We work with the remaining groups of
Section~\ref{sec:families}: unipotent series of general linear or unitary
groups, and the regular classical groups of
Lemma~\ref{fam:regular-datum}, with the sign packets of
Lemma~\ref{fam:packets}.  The letter $\mathbf H$ also allows their
successive split Levis.  Previously extracted linear factors can be
retained as fixed factors; their ordinary parameters and labels are
fixed throughout an extraction.  In the runner applications those
factors have order prime to $p$.

We first specify the ordinary character theory being used.  We then
prove the compatibility with ordinary Jordan labels that is needed
below.  Finally we show that the core projections in the runner
argument are projections by block idempotents.  All three statements
concern a fixed rational semisimple parameter in each Levi: we do not
sum over different Levi classes that fuse in the ambient group.

For a finite group $H$ and $\chi\in\Irr(H)$ put
\[
 \operatorname{def}_p(\chi)=\frac{|H|_p}{\chi(1)_p}.
\]
This is the degree defect, expressed as an order.  If a central
$p$-subgroup $Z$ is in the kernel of $\chi$, its degree defect as a
character of $H/Z$ is $\operatorname{def}_p(\chi)/|Z|$.

\begin{proposition}[Labels and ordinary branching]\label{prop:label-degrees}
The following conventions and formulas apply to the unipotent factors
of the ordinary Jordan labels in the groups just specified.
\begin{enumerate}
\item For $\GL_n^\epsilon(q)$, where $\epsilon=1$ or $-1$,
unipotent characters are indexed by partitions $\nu$ of $n$, and
\[
 \operatorname{def}_p(\chi_\nu)
   =\prod_{h\in\operatorname{Hook}(\nu)}
       \bigl((\epsilon q)^h-1\bigr)_p.
\]
Here the $p$-part of a negative integer means the $p$-part of its
absolute value.
\item In types $B,C,D$ a symbol is a pair $(A,B)$ of strictly
increasing finite sets of nonnegative integers.  Simultaneously
replacing both rows $A,B$ by $\{0\}\cup(A+1),\{0\}\cup(B+1)$
does not change the symbol.  The rows are unordered, subject to the
usual two separate labels for equal rows in split type $D$ of positive
rank.  If $u=|A|+|B|$, the rank is
\[
 \sum_{a\in A}a+\sum_{b\in B}b
       -\left\lfloor\frac{(u-1)^2}{4}\right\rfloor.
\]
The difference $|A|-|B|$ is odd in $B,C$, is $0$ modulo $4$ in
split $D$, and is $2$ modulo $4$ in nonsplit $D$.  We use one
empty-group label in rank zero, and the toral convention for
$\mathrm{SO}_2^\pm$.

A hook is a bead $a$ and a smaller gap $b$ in the same row; a cohook
uses a gap in the other row.  Their lengths are $a-b>0$.
For a packet with field parameter $q_i$, its degree defect at the
odd prime $p$ is the product of
$(q_i^h-1)_p$ over hooks and $(q_i^c+1)_p$ over cohooks.
The inherited central torus contribution must also be included.
For an original regular group this is its full connected center;
for a residual factor $\mathbf R$, it is the torus $\mathbf T_R$
of Lemma~\ref{fam:regular-datum}.  The product over packets,
including a toral $D_1$ packet, multiplied by $|\mathbf T_R^F|_p$,
is the degree defect of the ordinary Jordan correspondent.
Thus the toral packet is counted exactly once.
\item Lusztig induction from a single positive cycle torus of odd length
$d$ and a residual classical group adds $d$-hooks to symbols;
a negative cycle adds $d$-cohooks.  The coefficients have the usual
leg signs, including the common extra minus sign for a cohook
in type $D$.  Each removal contributes absolute value one before the
type-$D$ identifications.  The individual split labels are retained.
Removing all cycles of the relevant length packs the runners and
gives a core.  Every unipotent label with that core occurs with
nonzero coefficient in induction from the core and its cycle tori,
with the appropriate residual orthogonal sign and, when necessary,
the appropriate choice of split label.
\item Ordinary split Harish--Chandra restriction removes $1$-hooks
for $\GL$, $2$-hooks for $\GU$ with a
$\GL_1(q^2)$ extraction, and same-row $1$-hooks for symbols.
On an ordinary Harish--Chandra series this is Young branching on
partitions or bipartitions.  In the unitary case the $2$-core stays
fixed; for symbols the row-length difference stays fixed up to
exchange.  The relative Weyl group is $W(B_b)$, except for the
zero-difference series in type $D$, where it is $W(D_b)$ and the
branching keeps the two labels attached to equal bipartitions.

At the first sign wall over a cuspidal symbol, the ratio of the two
ordinary rank-one degrees is a power of $q_i$, up to inversion.
Over the unitary cuspidal staircase of size $t(t+1)/2$ the ratio is
$q^{2t+1}$, up to inversion.
\end{enumerate}
\end{proposition}

\begin{proof}
These are the ordinary partition and symbol parametrizations, degree
formulas, and Harish--Chandra theory of
Lusztig~\cite{LusztigCharacters}.  In the degree formulas the factors
that have been omitted are powers of the defining characteristic,
signs, and, for symbols, powers of $2$; none changes the asserted
odd $p$-part.  The degree rule for Jordan decomposition cancels the
ambient-to-centralizer index, leaving exactly the centralizer degree
defect.  Central isogenies preserve the rational order and unipotent
degrees, so these calculations apply to our regular forms and to
Frobenius orbits of factors, using $q_i=q^\delta$.

The single-cycle formulas are Asai's hook and cohook formulas
\cite{Asai}; see also \cite[Theorem~3.2]{LubeckMalle}, where a
degenerate symbol initially denotes the sum of its two characters.
The individual normalization needed here is spelled out in
Lemma~\ref{lab:sparse} below.  The assertion about a full packed core
is the ordinary unipotent cyclotomic Harish--Chandra series rule
\cite[Theorem~3.2 and Section~3.A]{BMM}.
This is a statement about ordinary unipotent characters.
For split Harish--Chandra series the ordinary endomorphism algebras
give the indicated Weyl-group branching.  The rank-one degree
ratios are the ordinary rank-one parameters, equivalently the
ratios in the same partition and symbol degree formulas.
\end{proof}

Here and below the $e$-part of a torus means its connected
$\Phi_e$-subtorus, defined by the $\Phi_e$-isotypic subspace of its
rational cocharacter space.  Taking the full lattice in that subspace
defines the actual subtorus.  Set
\[
 e=\ord_p(q),\qquad e_i=\ord_p(q_i),\qquad
 (d,\eta)=
 \begin{cases}
  (e_i,+1),&e_i\text{ odd},\\
  (e_i/2,-1),&e_i\text{ even}.
 \end{cases}
\]
The core operation on a packet removes positive $d$-cycles when
$\eta=+1$ and negative $d$-cycles when $\eta=-1$.
The ambient cyclotomic index remains $e$, including when
$q_i=q^2$.

\subsection{Uniform tests and a minimum principle}

For a fixed target character, the coefficients of an induction matrix
form an integral vector on the source characters.  Torus pairings
need not determine this vector.  For the unipotent single-cycle
formulas of Proposition~\ref{prop:label-degrees}, we will show that
the hook or cohook coefficient vector is the unique integral vector
of minimum norm with its prescribed pairings.  The transport proof
in the next subsection will then identify the ordinary induction
rows by comparing the sums of their squared norms.

A uniform function is a linear combination of Deligne--Lusztig torus
characters.  The torus almost characters provide convenient tests
for its scalar products with irreducible unipotent characters.
We record precisely the portion of the classical Fourier formula
that will be used.  Within a symbol family, the entries occurring
twice are fixed, as is the set $\Omega$ of entries occurring once.
Write $h=|\Omega|$.  A row of the Fourier matrix is indexed by a
permitted distribution $B\subseteq\Omega$ between the symbol rows.
For types $B,C$ we orient the symbol to have row-length difference
$1$ modulo $4$; the resulting subsets have one fixed parity.
For type $D$ the permitted differences are even, and complementary
subsets represent row exchange.  The uniform columns are balanced
distributions $U$, of cardinality $\lfloor h/2\rfloor$ or the
complementary cardinality.  In type $D$ complementary columns are
identified.  For graph-twisted $D$ the columns use extensions of the
nondegenerate, equal-row-length Weyl labels to the twisted coset.
Equal row lengths here refer to symbol defect, not to block defect.

Write $f_B(U)$ for the scalar product in this rectangle.  The
classical Fourier formula~\cite{LusztigCharacters}, in the
subset conventions of
\cite[Section~6.1 and Proposition~6.3]{DigneMichel1990}, says that
pairings between different families vanish, that all entries in
one rectangle have the same nonzero absolute value, and that
\begin{equation}\label{lab:fourier-ratio}
 \frac{f_B(U)f_{B'}(U')}{f_B(U')f_{B'}(U)}
   =(-1)^{|(B\mathbin\triangle B')\cap
                    (U\mathbin\triangle U')|}.
\end{equation}
Phases attached to entire test columns have no effect on this
identity.  Degenerate split symbols have singleton families and
are separately uniform, with their own degenerate Weyl tests.
Type $A$ is uniform.  Products of packets use tensor products of
these test matrices.

The classical formulas hold over every finite field by
\cite[Corollary~4.3]{AsaiSmall}.  The following separation assertion
is the classical case of Digne--Michel
\cite[Proposition~6.3]{DigneMichel1990}.
It concerns ordinary characters only and is independent of the
coefficient prime $p$.  In particular, it will also apply at $p=2$
in Section~\ref{sec:periodicity}.
We recall the subset argument to fix the conventions used below;
the marked-diagram consequence is then applied to our Levi data.

\begin{lemma}[Separation by uniform tests]\label{lab:fingerprints}
The uniform pairing vectors of two ordinary unipotent labels in
these groups are proportional only when the labels coincide.
Consequently ordinary Jordan labels are functorial under rational
conjugacies of marked torus and Levi diagrams.
\end{lemma}

\begin{proof}
Different families have disjoint nonzero test supports.  In one
nondegenerate family, the differences $U\triangle U'$ of all
balanced subsets span the even-weight subspace of $\F_2^\Omega$:
for any distinct $i,j$, choose a balanced subset containing $i$
and not $j$, and replace $i$ by $j$.  These differences generate
all $\{i,j\}$.  Thus proportionality and
\eqref{lab:fourier-ratio} imply that
$B\triangle B'$ is either empty or all of $\Omega$.
In types $B,C$, $h$ is odd and the fixed parity convention excludes
the latter possibility.  In type $D$ it is precisely row exchange.
The cases $h=1$, or $h=2$ in type $D$, have only one relevant row;
degenerate labels have their separate tests.  Tensor factors may be
varied separately, proving the product assertion.

The ordinary Jordan torus-pairing rule is invariant under transport
of the full marked dual torus diagram.  Such transport therefore
preserves every uniform pairing of the prospective label.  The
first assertion removes the only possible character permutation
ambiguity.
\end{proof}

\begin{lemma}[The separated slice]\label{lab:slice}
Fix a nondegenerate Fourier rectangle and two positions
$P=\{x,y\}\subset\Omega$.  Suppose an integral vector $v$ on its
rows has two nonzero entries, both of absolute value one, at the
distinct labels $B$ and $B\triangle P$.  Suppose its uniform
pairing vector vanishes when $|U\cap P|\ne1$ and has twice the
single-entry absolute value when $|U\cap P|=1$.
Among integral vectors with this pairing vector, $v$ is the unique
one of minimum squared Euclidean norm, namely two.
The same conclusion holds after tensoring with fixed rows on
other packets.
\end{lemma}

\begin{proof}
Let $c>0$ be the common absolute value in the rectangle and put
$\mathcal S=\{U:|U\cap P|=1\}$.  This set is nonempty.
A vector of squared norm at most one is zero or one signed unit
vector, so its pairings have absolute value at most $c$.
It cannot give the required value $2c$ on $\mathcal S$.
An integral vector of squared norm two has exactly two distinct
signed unit entries.  On every column in $\mathcal S$, equality
in the triangle inequality forces each of these entries to give
the same value as each summand of $v$.

We determine which signed row can agree with a given signed row
on $\mathcal S$.  Regard subsets as vectors over $\F_2$ and let
$E(P^c)$ denote the even-weight subspace supported on
$P^c=\Omega\setminus P$.  If $h\ge3$, then
\begin{equation}\label{lab:slice-span}
 \left\langle U\triangle U':U,U'\in\mathcal S\right\rangle
     =\langle\mathbf1_P\rangle\oplus E(P^c).
\end{equation}
Indeed one may switch the chosen element of $P$.  In $P^c$ the
chosen subsets have size one less than the balanced size.  For
$h\ge4$ that size lies strictly between zero and $|P^c|$, so the
same single exchange argument as above generates $E(P^c)$.
For $h=3$ the complement has one element and its even subspace
is zero, giving the same formula.  In type $D$, lift balanced
columns to both complementary representatives before making this
calculation.  Even row differences ensure that their pairing
ratios agree on the two representatives.

The annihilator of the space in \eqref{lab:slice-span} consists
of vectors constant on $P$ and constant on $P^c$.  In odd type,
the fixed row parity excludes toggling $P^c$, whose size is odd.
In type $D$, toggling $P^c$ is the same as toggling $P$ modulo
full complement.  Thus the only possible rows are $B$ and
$B\triangle P$.  Their signs are fixed by a single column in
$\mathcal S$; the two distinct entries are consequently exactly
those of $v$.  The case of two distinct rows does not occur for
$h\le2$.  Varying test columns independently in every other
packet and applying Lemma~\ref{lab:fingerprints} fixes the inactive
rows and proves the tensor-product assertion.
\end{proof}

\begin{lemma}[Sparse single-cycle rows]\label{lab:sparse}
Fix a target unipotent label and a source family for one of the
single-cycle inductions of Proposition~\ref{prop:label-degrees}.
The list of coefficients on that family is zero, one signed unit
entry, or two distinct signed unit entries.  In the last case it
satisfies Lemma~\ref{lab:slice}, with $P$ the two moved positions.
In all cases this list is the unique integral list of minimum
squared norm having its uniform pairings.
\end{lemma}

\begin{proof}
Pad the two beta rows by simultaneous shifts so that the row-length
conventions agree in all symbols being compared.  The two family
multisets determine the starting and finishing positions
$a,a-d$ of a removal.  Away from split degeneracy there is at
most one removal in each of the two rows.  Two removals occur
only when the target has two beads at $a$ and no bead at $a-d$.
Their source distributions differ by toggling
$P=\{a,a-d\}$.  If these two removals gave the same unordered
source label, the target rows would agree away from these two
positions and also at them; that is exactly a degenerate target.
Thus, in the present case, the two source labels are distinct.

We verify their uniform support; this also fixes their relative
sign.  On Weyl-group tests, restriction at either a positive or
a negative signed cycle is same-row hook removal on bipartitions.
This follows from the induced-character description of signed
permutation characters and the symmetric-group
Murnaghan--Nakayama rule: the sign of the cycle changes the sign
of one bipartition contribution, but does not change which row
loses the hook.  A balanced source distribution can reach the
target duplicate at $a$ and vacancy at $a-d$ only when it
separates those two positions.  Hence torus transitivity forces
the two-entry pairing to vanish off $\mathcal S$.
Formula~\eqref{lab:fourier-ratio} says that the two row ratios are
constant on each slice and opposite on the two slices.  The two
rows are not proportional by Lemma~\ref{lab:fingerprints}.
Consequently their signed sum has absolute value $2c$ on
$\mathcal S$.  Both slices occur in every relevant two-removal
family.  Lemma~\ref{lab:slice} applies.

Here is the normalization at the type-$D$ boundary.  A degenerate
active character is uniform, and a contributing family on the
other side has two singles: deleting one bead from equal rows,
and inserting it at a previously vacant position, leaves just
these two single positions.  There is one ordinary row and one
uniform column in that nondegenerate rectangle.  The relevant
cycle classes do not split under restriction from $W(B)$ to
$W(D)$.  A positive cycle here has odd length, so flipping all
its coordinate signs is an odd-sign element of its centralizer;
a negative cycle itself supplies such an element.  Thus either
half of an equal-bipartition Weyl character has half the $W(B)$
value on the cycle classes in question.  Its two row removals
combine to give absolute value at most one on the residual
unordered Weyl label.  Conversely, from a fixed orientation of
a nondegenerate bipartition, only one hook reaches the equal
rows, and its coefficient on each split label has absolute value
one.  On the twisted coset the same calculation uses the
extensions of the nondegenerate Weyl characters; switching the
extension changes an overall sign, not the absolute value.
It therefore never identifies the two split labels.

For completeness, these checks also cover the ends of the rank
range.  A toral $D$ packet has its single unipotent character.
For a rank-zero residual packet, evaluate the Weyl character at
the full signed cycle.  In type $D$ only the single-hook
bipartitions just considered contribute.  In types $B,C$ their
symbol families have one or three singles.  In the three-single
case the full hook can leave either packed empty row; there are
at most two balanced contributions, each of Fourier absolute
value $1/2$, so again an individual coefficient has absolute
value at most one.  The one-single case has one test.  These
computations are also the individual-character interpretation
of the hook/cohook formulas when the degenerate-symbol notation
in \cite[Theorem~3.2]{LubeckMalle} is used.

Finally, a zero list is the unique vector of norm zero.  For one
signed entry, any integral competitor of norm at most one must
itself be one signed entry, and Lemma~\ref{lab:fingerprints}
identifies it.  Inactive factors are fixed and tensor with the
active calculation.  This proves the assertion in every case.
\end{proof}

\subsection{Transport of a single cycle}

The minimum-norm statements now reduce single-cycle transport to
an equality of total norms.  We prove that equality by Mackey's
formula, retaining the specified rational parameter.
We adapt the norm-comparison and minimum-norm argument of
Enguehard~\cite[Section~5.3, equations~(5.3.3)--(5.3.8) and
Lemma~5.3.9]{EnguehardJordan}.  We give the single-cycle argument
explicitly, retaining the rational Levi parameter, individual split
labels, toral and empty packets, and fixed inactive factors.

Let $s\in\mathbf H^{*F}$ be the fixed semisimple parameter and
put $\mathbf C=C_{\mathbf H^*}(s)$.  This centralizer is connected.
We allow any fixed additional linear packets, including parameters
with nontrivial $p$-parts on those inactive packets.  The active
packet is one of the classical packets just described.
For this subsection we also allow ambient Levis obtained by earlier
removals of disjoint cycles of the indicated kind.  The removed
cycles are fixed torus packets in their semisimple centralizers;
on the ambient side they may lie in linear factors.  These groups
are still Levis of the original regular group, so both centers
remain connected and both derived groups remain simply connected.
Their active centralizer is the smaller residual classical packet.
Thus this enlarged class is closed under removing another disjoint
cycle and under the intersection diagrams used in the proof below.

In a rational maximal torus of $\mathbf C$ choose $\mathbf V$
supported on one cycle only: either a split hyperbolic coordinate,
or the primitive $\Phi_e$-part of one of the signed $d$-cycles.
Put
\begin{equation}\label{lab:cycle-levi}
 \mathbf M^*=C_{\mathbf H^*}(\mathbf V),\qquad
 \mathbf C_M=C_{\mathbf M^*}(s)=C_{\mathbf C}(\mathbf V).
\end{equation}
The second group is, up to central isogeny, the residual classical
factor times the cycle torus and the inactive factors.  To see this
on the root datum,
in each absolute component the nonzero coordinate characters on
$\mathbf V$ are distinct up to sign.  On the primitive part they
are the successive powers of a primitive $e_i$th root, with the
prescribed signed closing.  All other coordinates vanish on
$\mathbf V$.  Precisely the desired residual roots vanish.
For a packet with $q_i=q^\delta$, its full signed orbit has length
$\delta d$ before folding.  The split-coordinate construction
gives an ordinary split Levi and a split parabolic.

Write $J_{H,s}$ and $J_{M,s}$ for the ordinary Jordan bijections
to the unipotent characters of $\mathbf C^F$ and
$\mathbf C_M^F$.  Let $\varepsilon_{\mathbf K}$ denote the usual
split-rank sign of an $F$-group $\mathbf K$.

\begin{proposition}[Single-cycle Jordan transport]\label{prop:single-cycle}
With the fixed markings in \eqref{lab:cycle-levi}, the matrix of
$R^{\mathbf H}_{\mathbf M}$ on
$\mathcal E(\mathbf M^F,s)$, in ordinary Jordan labels, is
\begin{equation}\label{lab:single-cycle-matrix}
 A=\epsilon B,\qquad
 B=\left[R^{\mathbf C}_{\mathbf C_M}\right]_{\mathrm{unip}},
 \qquad
 \epsilon=
 \varepsilon_{\mathbf H}\varepsilon_{\mathbf M}
 \varepsilon_{\mathbf C}\varepsilon_{\mathbf C_M}.
\end{equation}
For the split hyperbolic extraction the sign is positive.
This includes individual split labels, toral and empty residual
packets, and products with fixed inactive packets.  Successive
split extractions consequently give ordinary partition,
bipartition, and folded type-$D$ Harish--Chandra branching on
these Jordan labels.
\end{proposition}

\begin{proof}
All matrix entries in $A$ are integers, since Lusztig induction
is the alternating character of a cohomology complex.  Series
disjunction puts its image in the stated ambient series.  On
uniform inputs, the equality \eqref{lab:single-cycle-matrix}
already follows from transitivity of torus induction and the
ordinary Jordan torus-pairing rule.  More explicitly, apply
transitivity to each torus character in $\mathbf M$ with
parameter $s$ and pair with a fixed target character.  The
Jordan torus pairings on the two sides give its predicted
uniform pairing against every source family.  By
Lemma~\ref{lab:sparse}, that target row of $\epsilon B$ has
the unique minimum norm among integral rows with these pairings.
Summing over source families and then target rows gives
\begin{equation}\label{lab:norm-min}
 \|A\|_{\mathrm{HS}}^2\ge\|B\|_{\mathrm{HS}}^2,
\end{equation}
and equality forces $A=\epsilon B$.  Here the square of the
Hilbert--Schmidt norm is the sum of the squares of all matrix
entries.  It remains to prove equality of these total norms.

\smallskip\noindent\emph{Matching the rational terms.}
We use the Lusztig-induction Mackey formula, which holds for
these classical components for every $q$
\cite[Theorem~3.8(4)]{BonnafeMichel}.
The argument is by induction on the rank of the active packet.
Apply Mackey to
$({R^{\mathbf H}_{\mathbf M}})^*R^{\mathbf H}_{\mathbf M}$
and take its trace on $\mathcal E(\mathbf M^F,s)$.
We describe the rational terms and their markings, since an
unmarked Weyl-double-coset count would not suffice.
The object to match is an intersection together with its two
rational embeddings into the Levi.  Both embeddings are needed
to impose the chosen $s$-series at both ends of each term.

Fix dual $F$-stable maximal tori $\mathbf T\subseteq\mathbf M$
and $\mathbf T^*\subseteq\mathbf M^*$, and retain their full
character and cocharacter lattices.  Put
$W=W(\mathbf H,\mathbf T)$ and $W_M=W(\mathbf M,\mathbf T)$.
The algebraic double positions in the Mackey formula are indexed
by $W_M\backslash W/W_M$.  If $n\in N_{\mathbf H}(\mathbf T)$
represents $w\in W$, their
intersection is
$\mathbf K=\mathbf M\cap{}^n\mathbf M$; it is the centralizer
of the torus generated by $Z^\circ(\mathbf M)$ and
${}^nZ^\circ(\mathbf M)$, and hence is connected.  Its roots
are $\Phi_M\cap w\Phi_M$.  The stabilizer for the double action
of $\mathbf M\times\mathbf M$ is the graph of this connected
intersection.  Lang's theorem therefore gives exactly one
rational double orbit for each $F$-stable algebraic double orbit.
This assertion does not identify the different rational conjugacy
classes of maximal tori in an intersection.

We record the descent of both embeddings, since an $F$-stable
Weyl double coset need not have an $F$-fixed representative.
Choose $u,v\in W_M$ with $F(w)=uwv^{-1}$, lift $u$ to
$a\in N_{\mathbf M}(\mathbf T)$, and set
$b=F(n)^{-1}an\in N_{\mathbf M}(\mathbf T)$.
Thus $F(n)=anb^{-1}$ and $b$ lifts $v$.  The maps
\[
 j_1:\mathbf K\hookrightarrow\mathbf M,\qquad
 j_2=\operatorname{Ad}(n^{-1}):\mathbf K\longrightarrow\mathbf M
\]
commute with the respective Frobenius maps
\[
 F_1=\operatorname{Ad}(a^{-1})F,\qquad
 F_2=\operatorname{Ad}(b^{-1})F,\qquad F_K=F_1|_{\mathbf K}.
\]
Choose $x,y\in\mathbf M$ with $F(x)=xa^{-1}$ and
$F(y)=yb^{-1}$.  Then $g=xny^{-1}\in\mathbf H^F$.
Conjugating $\mathbf K$ by $x$ gives
$\mathbf M\cap{}^g\mathbf M$ with its two rational embeddings,
inclusion and $\operatorname{Ad}(g^{-1})$.

Here is the dual Frobenius check on the full lattices.
Choose dual bases of $X(\mathbf T)$ and
$X(\mathbf T^*)=Y(\mathbf T)$.
Write $Q$ for the pullback of $F$ on $X=X(\mathbf T)$, with
Weyl elements acting on $X$ in the usual root-datum convention.  Then
\[
 F(w)=Q^{-1}wQ,\qquad F_1^*=Qu,\quad F_2^*=Qv,
 \qquad wQv=Quw.
\]
Write $F^\vee$ for the Frobenius on the dual group, also
denoted by $F$ elsewhere.  On its character lattice
$Y=X(\mathbf T^*)$ put
$Q^\vee=Q^t$ and $w^\vee=w^{-t}$.  Transposing gives
\[
 w^\vee v^tQ^\vee=u^tQ^\vee w^\vee.
\]
Consequently, for
$U=(Q^\vee)^{-1}u^tQ^\vee$ and
$V=(Q^\vee)^{-1}v^tQ^\vee$,
\[
 F^\vee(w^\vee)=Uw^\vee V^{-1}.
\]
Apply the same normalizer-lift and Lang construction on the dual
side.  Its two twisted Frobenius pullbacks are
$Q^\vee U=(Qu)^t$ and $Q^\vee V=(Qv)^t$, and its second
embedding has pullback $w^\vee$.  Its intersection roots are
$\Phi_M^\vee\cap w^\vee\Phi_M^\vee$.
Thus the construction dualizes both embeddings and their
Frobenius maps, including the central lattices.

For completeness, fix $n$ and compare two choices $(a,b)$ and
$(a',b')$.  There is $k\in N_{\mathbf K}(\mathbf T)$ with
\[
 a'=ak,\qquad b'=b(n^{-1}kn).
\]
The new intersection Frobenius is
$F'_K=\operatorname{Ad}(k^{-1})F_K$.  Choose
$z\in\mathbf K$ with $F_K(z)=zk^{-1}$.
Conjugation by $z$ on the first endpoint and by $n^{-1}zn$
on the second identifies the two twisted diagrams, and hence
their descents.  Indeed, one may take $x'=xz$ and
$y'=y(n^{-1}zn)$, which gives $x'ny'^{-1}=xny^{-1}$.
Replacing
$n$ by $lnr^{-1}$ for $l,r\in N_{\mathbf M}(\mathbf T)$
changes only the two endpoint markings: take
$a'=F(l)al^{-1}$ and $b'=F(r)br^{-1}$.
Different Lang solutions $x,y$ differ on the left by elements
of $\mathbf M^F$, so the resulting rational double orbit is
unchanged.  The same argument applies in the connected dual
intersection.  Lattice duality makes the construction reversible.
We have therefore matched the rational Levi-intersection diagrams
with both embeddings.  The auxiliary torus used for the lattice
calculation is discarded after descent; no rational conjugacy
of arbitrary choices of common maximal torus is asserted.

For a dual rational representative $g\in\mathbf H^{*F}$,
refine its term by rational semisimple classes in the intersection
that map to the class of $s$ on both sides.
Series disjunction is exactly this refinement of the projected
primal Mackey term.  Every such matching class can be written
\[
 {}^{m_1}s={} ^{g m_2}s,
       \qquad m_1,m_2\in\mathbf M^{*F}.
\]
Changing the two markings gives
$h=m_1^{-1}g m_2\in\mathbf C^F$.
The remaining changes are the left and right actions of
$\mathbf C_M^F$.  Conversely a double coset with a
representative $h\in\mathbf C^F$ gives this matching class.
Changing $g$ while returning to the same double position
conjugates the parameter by the intersection, exactly as in
the refinement.  Thus the refined terms are in bijection with
the Mackey terms for $(\mathbf C,\mathbf C_M)$.
All centralizers involved are connected, since the relevant
derived groups on the dual side are simply connected.
Lemma~\ref{lab:fingerprints}, applied to the transported torus
diagrams, identifies the Jordan labels under every conjugation
in this correspondence.

\smallskip\noindent\emph{Comparing the contributions.}
The rational Mackey terms and their parameter markings have now
been matched.  We next compare their contributions: equal supports
give the identity terms, whereas disjoint supports reduce to a
single-cycle calculation on a smaller active packet.
There is an explicit description of the intersections in this
bijection.  For $h\in\mathbf C^F$, the intersection of
$C_{\mathbf H^*}(\mathbf V)$ and
$C_{\mathbf H^*}({}^h\mathbf V)$ contains a maximal torus
if and only if $\mathbf V$ and ${}^h\mathbf V$ commute.
For the forward implication, both tori belong to the connected
centers of their respective Levis and hence to every maximal
torus in them.  The converse follows by placing the commuting
tori in a maximal torus.  The same condition in $\mathbf C$
gives its Mackey intersection
$C_{\mathbf C}(\mathbf V,{}^h\mathbf V)$.

Take a rational common maximal torus in this connected
intersection.  The full signed-cycle torus containing
$\mathbf V$ belongs to the connected center of
$C_{\mathbf C}(\mathbf V)$: within each absolute component
the restricted coordinate characters are nonzero and distinct
up to sign, so no root involving a cycle coordinate vanishes
except the roots of the inactive factors.  Every maximal
torus of this centralizer therefore contains that full cycle
torus.  Its active support is consequently preserved when
we pass to the common maximal torus, and the same holds for
${}^h\mathbf V$.

On a packet of $\delta$ components, $F$ permutes the
components transitively and $F^\delta$ cyclically permutes
the distinct within-component restricted characters with
the prescribed signed closing.  The support is one signed
$F$-orbit, even though restrictions on different components
can coincide.  The two supports are therefore equal or
disjoint.  For a cyclotomic cycle the characteristic
polynomial on the full orbit is $X^{\delta d}-\eta$,
so $\Phi_e$ occurs with multiplicity one.  For a split
extraction it is $X^\delta-1$, with $\Phi_1$ occurring
once.  Thus the indicated primitive or split subspace
is unique on its support.  Equal supports thus
mean equal subtori.  Conjugation in $\mathbf C$ cannot add
central coordinates or move a support to a different
eigenvalue packet.

If the subtori are equal, the maps in the corresponding
intersection term are identities followed by the matched
conjugation.  If they are disjoint, both induction maps in
that term remove a single cycle of the same kind from a
strictly smaller residual packet.  The cycle already removed
is now an inactive torus factor.  This remains true if
coincident coordinates from distinct sign packets formed
an ambient linear factor in $\mathbf M^*$: its intersection
with $\mathbf C$ is the indicated cycle torus.  By induction
on active rank, both maps in the term agree in Jordan labels
with the centralizer maps.  To check the signs, write
$\mathbf K^*=C_{\mathbf H^*}(\mathbf V,{}^h\mathbf V)$
and $\mathbf D=C_{\mathbf C}(\mathbf V,{}^h\mathbf V)$.
Each smaller map has relative sign
$\varepsilon_{\mathbf M}\varepsilon_{\mathbf K}
 \varepsilon_{\mathbf C_M}\varepsilon_{\mathbf D}$;
conjugating the marking preserves these split ranks.
The two signs therefore multiply to one.  The equal-support terms start
the induction, including rank zero and the toral $D$ case;
an ambient identity extraction needs no induction.

The trace of the projected ambient Mackey formula is
therefore the trace of the unipotent Mackey formula for
$(\mathbf C,\mathbf C_M)$.  These traces are respectively
$\|A\|_{\mathrm{HS}}^2$ and $\|B\|_{\mathrm{HS}}^2$.
Equality in \eqref{lab:norm-min} proves
\eqref{lab:single-cycle-matrix}.  For split extraction both
sides are ordinary Harish--Chandra induction matrices, with
nonnegative entries and a nonzero entry, so the relative sign
is positive.  Transitivity proves the last assertion.
\end{proof}

\begin{corollary}[Projected split induction]\label{lab:projected-hc}
Let $s$ be a $p'$-parameter and fix its rational class in a
split Levi $\mathbf M$ as above.  On characters in
$\mathcal E_p(\mathbf M^F,s)$, the coordinates in
$\mathcal E(\mathbf H^F,s)$ after split induction depend only
on the coordinates in $\mathcal E(\mathbf M^F,s)$.
On these coordinates the matrix is the nonnegative branching
matrix of Proposition~\ref{prop:single-cycle}, iterated when
several size-one factors are extracted.  The assertion remains
valid after a common central $p$-quotient.
\end{corollary}

\begin{proof}
A row with parameter $st$, where $t\ne1$ is a commuting
$p$-element, can induce only to that geometric parameter.
It cannot induce to the $p'$-parameter $s$, because its
$p$-part stays nontrivial under conjugacy.  Thus its projection
onto $\mathcal E(\mathbf H^F,s)$ is zero.  For $t=1$ use
Proposition~\ref{prop:single-cycle} and transitivity, always
retaining the chosen Levi class.  The assertion for a common
central quotient follows because its left and right actions
on the induction variety agree.  Ordinary split induction
and restriction are exact and preserve projectives, so the
same coordinate rule applies to projective characters and
to their block summands.
\end{proof}

\subsection{Core collections are block collections}

We have determined the ordinary matrices, including the
nonuniform rows.  To use them on projective modules, the core
projections must also be realized by block idempotents.  This is
the separate modular assertion proved in this subsection.

We use a restricted form of the good-prime block theorem of
Cabanes--Enguehard.  The relevant statements are
\cite[Theorems 2.5 and 4.1]{CE99}.
For the connected reductive groups at hand, ordinary
$e$-cuspidal pairs in $p'$-series have assigned blocks; all
constituents of their Lusztig induction lie in the assigned
block, and rationally nonconjugate pairs give different blocks.
An $e$-split Levi is a centralizer of a $\Phi_e$-torus; ordinary
$e$-cuspidality means vanishing of adjoint Lusztig induction to
every proper $e$-split Levi.  This is the general connected
reductive statement of the theorem; a central torus need not be
split.  In our application the semisimple components are
classical of type $B,C,D$ or $A_1$, $p$ is odd and good, both
ambient centers are connected, and the relevant center and
fundamental-group orders have no $p$-part.  There is no
triality component.  This includes $p=3$; see also the
classical small-prime discussion in \cite[Remark~5.2]{CE99}
and \cite[proof of Theorem~3.14(e)]{KessarMalle2015}.
We apply this theorem directly to these connected groups.
For the connected reductive form used here, including its possibly
nonsplit central torus, see also
\cite[Assumption~2.1.2 and Propositions~2.1.6--2.1.7]{EnguehardJordan}.

\begin{lemma}[The packed inducing pair]\label{lab:packed-pair}
Fix target Jordan labels on the active classical packets,
and let $\kappa$ be their collection of packed cores.
Take the signed cycles removed in packing these labels
to $\kappa$; their number in each packet is determined by
its target rank and core rank.  Centralize their independent
primitive $\Phi_e$-subtori in $\mathbf H^*$, obtaining
$\mathbf L^*$.  Then $s$ belongs
to $\mathbf L^{*F}$, and its dual Levi $\mathbf L$ is
$e$-split.  The ordinary character $\lambda$ of
$\mathbf L^F$ whose Jordan label is $\kappa$, trivial on
the removed cycle tori and with the fixed inert torus data,
has central $p$-defect and is ordinary $e$-cuspidal.
If there are separate split core labels the assertion applies
to the corresponding choices individually.
\end{lemma}

\begin{proof}
Independent subtori are obtained by giving each of these removed
signed cycles its own coordinates and making its cocharacters
zero elsewhere.  Their centralizer is an $e$-split Levi and
contains $s$.  The permitted torus positions and the residual
orthogonal sign are exactly those of the unipotent core
rule; a negative cycle switches the type-$D$ defect congruence.
That rule also supplies the inducing label when the residual
rank is zero.  We check the degree and cuspidality assertions
on this actual Levi.

A packed core has no hook or cohook divisible by the relevant
cyclotomic step, so its degree defect contributes no $p$-part
beyond its torus data.  The $p$-part of every removed cycle
torus is already contained in its primitive $\Phi_e$-part.
Indeed a factor $\Phi_j(q)$ can be divisible by $p$ only for
$j=e p^a$ with $a\ge0$ (with the usual $e=1$ convention).
The signed absolute cycle length is $\delta d$, with
$\delta\le2$ and $d\mid(p-1)$ or $2d\mid(p-1)$.
It has no $p$-factor, so $a>0$ cannot occur.  These primitive
tori and the ambient central torus are central in
$\mathbf L^*$ by construction.  No residual core torus
has an additional noncentral $\Phi_e$-part: such a toral
$D$ remainder would still admit a removable cycle.  Frozen
extracts have order prime to $p$.

Explicitly, a split $D_1$ packet has order $q_i-1$.
If its $p$-part is nontrivial, then $e_i=1$, and its label
admits a removable $1$-hook.  A nonsplit $D_1$ packet has
order $q_i+1$; a nontrivial $p$-part forces $e_i=2$ and
a removable $1$-cohook.  Neither can remain in a packed
core.  Thus every residual $D_1$ core torus has order
prime to $p$, even when it enlarges the full residual
center.  The $p$-part of the full center in
\eqref{lab:central-defect} is consequently the inherited
central contribution together with the removed primitive-cycle
torus contributions.

There is no odd-prime isogeny loss in this calculation.
Within each absolute packet the primitive-cycle coordinate
characters are distinct up to sign.  Equalities between
different sign packets can therefore tie at most two
coordinates and create at most $A_1$ systems.  The residual
semisimple systems are $B,C,D$, and frozen extractions add
only the stated rank-one systems.  Their standard center
orders are powers of $2$.  The signed equality lattices and
the regular central dressing likewise have at worst
$2$-primary indices on both root-data sides.  Thus rational
torus orders, computed on the rational spaces, give the
correct $p$-parts also for the actual centers.  Dual centers
have the same rational order.  The degree rule of
Proposition~\ref{prop:label-degrees} consequently gives
\begin{equation}\label{lab:central-defect}
 \operatorname{def}_p(\lambda)=|Z(\mathbf L)^F|_p.
\end{equation}
One may also see the lower bound in this equality directly:
a character in a $p'$-series is trivial on the central
$p$-subgroup.

We have established the actual-center degree identity.
To prove ordinary $e$-cuspidality, we now combine projectivity
modulo that center with the larger central $p$-subgroup in every
proper $e$-split Levi.
Put $Z=O_p(Z(\mathbf L)^F)$.  Equation
\eqref{lab:central-defect} says that $\lambda$, viewed on
$\mathbf L^F/Z$, is an ordinary defect-zero character,
hence the character of a projective module.  Let
$\mathbf K$ be a proper $e$-split Levi of $\mathbf L$.
Its adjoint Lusztig restriction
$\psi={}^*R^{\mathbf L}_{\mathbf K}\lambda$ is a virtual
projective character of $\mathbf K^F/Z$.
Here is the quotient justification.  In a Deligne--Lusztig
induction variety the left and right actions of the common
finite central subgroup agree, and these translations are
free.  After division by $Z$, stabilizers for either group
action have prime-to-$p$ order: they come from unipotent
stabilizers, with the harmless prime-to-$p$ central kernels.
The projective-complex theorem for Rickard cohomology
\cite{RickardEtale} therefore sends projectives to virtual
projectives.  Over characteristic zero, taking the quotient
variety selects precisely the $Z$-invariant summand.  This
proves the assertion about $\psi$ without imposing modular
compatibility on Jordan decomposition.

Every parameter occurring in $\psi$ is conjugate to the
$p'$-element $s$.  Thus $\psi$ is trivial on
$O_p(Z(\mathbf K)^F)$.  This group is strictly larger than
$Z$.  Indeed the centers of these Levis are connected, and
proper $e$-split centralization adds a nonzero
$\Phi_e$-subspace to the center.  The quotient of the
centers is a torus whose order has a factor
$\Phi_e(q)$.  Lang's theorem on the connected kernel
$Z(\mathbf L)$ makes the sequence on rational centers
surjective, so the quotient adds a nontrivial $p$-part.
Choose a central element of order $p$ in its
image in $\mathbf K^F/Z$.

A virtual projective character vanishes on $p$-singular
elements.  For any $p$-regular element $x$ of
$\mathbf K^F/Z$, multiplying $x$ by this central element
gives a $p$-singular element, with the same value of $\psi$
because the central element is in its kernel.  Hence
$\psi(x)=0$ as well, and $\psi=0$.  Since $\mathbf K$
was arbitrary, $\lambda$ is ordinary $e$-cuspidal.
\end{proof}

\begin{remark}\label{lab:jordan-cuspidal}
The same pairs also have the centralizer description often
called $e$-Jordan cuspidality.  The $\Phi_e$-center of
$C_{\mathbf L^*}(s)$ is central in $\mathbf L^*$, and the
unipotent core is cuspidal for the corresponding cyclotomic
index.  If a packet is represented over $q^2$, projection of
an ambient $\Phi_e$-subtorus gives its cyclotomic subtorus
for $\ord_p(q^2)$; thus folding does not change which
cuspidality test is made.  Moreover
$\mathbf L^F/(Z(\mathbf L)^F[\mathbf L,\mathbf L]^F)$
and the intersection in this central product have
prime-to-$p$ order.  Clifford restriction across this
abelian quotient preserves degree $p$-parts: inertia
multiplicities are degrees of projective representations
of a $p'$-group, realized over a finite $p'$-central
extension.  Equation~\eqref{lab:central-defect} therefore
gives defect zero on the derived group after the central
part is removed, which is the quasi-central defect
condition.  The direct ordinary cuspidality proof above
is what is needed for the stated theorem of Cabanes--Enguehard.
\end{remark}

\begin{proposition}[Core block projections]\label{prop:core-blocks}
In each remaining type-$A$ unipotent target, and in each
regular classical target with fixed $p'$-parameter $s$,
the following holds.  Partition the basic rows in
$\mathcal E(\mathbf H^F,s)$ by their packed core
collections, using the actual unordered-symbol and
type-$D$ identifications.  Each part is the intersection
of that basic set with a union of blocks in
$\mathcal E_p(\mathbf H^F,s)$.  Distinct core collections
cannot occur in the same block.  Split sign choices
within one underlying core may be merged.

This remains true in every intermediate split Levi with
the specified frozen extractions, and after the common
central $p$-quotient.  Consequently all these core
projections are projections by sums of block idempotents
on projective modules.
\end{proposition}

\begin{proof}
For the general linear and unitary unipotent targets this
is the ordinary block/core rule of
Fong--Srinivasan~\cite{FongSrinivasan}, with partition step
$\ord_p(\epsilon q)$.

Consider a regular classical target.  Choose a basic row
with a given core collection and apply
Lemma~\ref{lab:packed-pair} to its labels.
The resulting pair is ordinary $e$-cuspidal, belongs to a
$p'$-series, and has central defect.  The preceding
root-lattice calculation verifies all the stated
classical hypotheses of the Cabanes--Enguehard theorem
on $\mathbf H$ and on the intermediate Levis: no
exceptional or triality component occurs, all relevant
odd primes are good, both centers are connected, and
the extra equality components are at most $A_1$.
Thus the theorem assigns a block to this pair.
By Proposition~\ref{prop:single-cycle}, transitivity,
and the ordinary unipotent core rule, every basic row
with the given core occurs in the induction of one
of its packed inducing characters.  Constituent
inclusion puts it in the corresponding assigned block.

We must distinguish the pairs for different cores
under rational conjugacy.  Suppose two such primal
pairs are rationally conjugate.  Transport a rational
torus in their Levis, with its full ambient marking.
Duality and the descent of the marked diagram used in the
proof of Proposition~\ref{prop:single-cycle} transport
the dual Levi diagram, and the series rule transports
its parameter.  Adjusting inside the target dual Levi
then makes this parameter exactly $s$.  The resulting
transport lies in $\mathbf C^F$.  Its action on the
residual unipotent labels is determined by the torus
pairings, hence by Lemma~\ref{lab:fingerprints}.

Inside $\mathbf C^F$ these are the ordinary
cycle-Levi and core data on the fixed eigenvalue
packets.  The central primitive-cycle factors record
the number of removed cycles in each nontoral packet;
a residual toral core has no additional noncentral
$\Phi_e$-part.  A conjugation matching the cycle data
acts on the residual packed label by its actual
symbol identifications.  It cannot change it to a
different row pair, nor can inner conjugation in the
connected centralizer exchange different eigenvalue
packets.  A wholly toral packet has only its single
unipotent label.  The allowed type-$D$ graph and sign
identifications give exactly the already stated
core equivalence.  Equivalently, this is the rational
nonconjugacy assertion for the packed data in the
ordinary unipotent core rule.  Therefore different
core collections give nonconjugate inducing pairs.
Uniqueness up to conjugacy in the Cabanes--Enguehard theorem puts
their rows in different blocks.

Every block in $\mathcal E_p(\mathbf H^F,s)$ meets
$\mathcal E(\mathbf H^F,s)$, and that intersection
is an integral basic set for the block.  Since the
basic rows have just been partitioned by disjoint
assigned block collections, each core collection
is exactly the basic-row intersection of its union
of blocks.  No assertion about modular Jordan
decomposition is involved.  The argument applies
verbatim with the frozen factors and their markings.
Finally a central $p$-quotient gives the usual
bijection of blocks, preserving their basic rows
with trivial central action.  Sums of the resulting
block idempotents preserve projectives, proving
the last assertion.
\end{proof}
